\documentclass[11pt,reqno]{amsart}

\usepackage[T1]{fontenc}
\usepackage{lmodern}
\usepackage{microtype}
\usepackage{amsmath,amssymb,mathtools}
\usepackage{mathrsfs}
\usepackage{enumitem}
\usepackage{booktabs}
\usepackage{xcolor}
\usepackage[colorlinks=true,
  linkcolor=blue!55!black,
  citecolor=green!40!black,
  urlcolor=blue!60!black,
  pdftitle={A Centered Volterra--Fredholm Realization of the Riemann Xi-Function},
  pdfauthor={Jingwei Liu},
  pdfsubject={Regularized Fredholm determinant and operator realization of the Riemann Xi-function},
  pdfkeywords={Riemann Xi function, Fredholm determinant, Volterra operator, Hilbert-Schmidt operator, Moore-Penrose inverse, PT symmetry}]{hyperref}
\usepackage[nameinlink,noabbrev]{cleveref}

\allowdisplaybreaks
\numberwithin{equation}{section}

\newtheorem{theorem}{Theorem}[section]
\newtheorem{proposition}[theorem]{Proposition}
\newtheorem{lemma}[theorem]{Lemma}
\newtheorem{corollary}[theorem]{Corollary}
\newtheorem{definition}[theorem]{Definition}
\newtheorem{remark}[theorem]{Remark}
\newtheorem{example}[theorem]{Example}
\newtheorem{problem}[theorem]{Problem}

\newcommand{\C}{\mathbb C}
\newcommand{\R}{\mathbb R}
\newcommand{\N}{\mathbb N}
\newcommand{\1}{\mathbf 1}
\newcommand{\HS}{\mathcal S_2}
\newcommand{\BTheta}{\mathcal B_{\Theta}}
\newcommand{\Vp}{\mathfrak V_p}
\newcommand{\Dom}{\operatorname{Dom}}
\newcommand{\Ran}{\operatorname{Ran}}
\newcommand{\Ker}{\operatorname{Ker}}
\newcommand{\tr}{\operatorname{tr}}
\newcommand{\Spec}{\operatorname{Spec}}
\newcommand{\sgn}{\operatorname{sgn}}
\newcommand{\dd}{\,\mathrm d}
\newcommand{\detTwo}{\det\nolimits_{2}}

\title[A centered Volterra realization of the Riemann $\Xi$-function]
{A Centered Volterra--Fredholm Realization\\
of the Riemann $\Xi$-Function}

\author{Jingwei Liu}
\date{August 28, 2026}

\subjclass[2020]{Primary 47G10, 47B10; Secondary 11M26, 47A10, 60E10,
81Q60}
\keywords{Riemann Xi function, regularized Fredholm determinant,
Volterra operator, Hilbert--Schmidt operator, Moore--Penrose inverse,
PT symmetry, supersymmetric quantum mechanics, pseudo-Hermitian operator}

\begin{document}
\raggedbottom

\begin{abstract}
Within the Hilbert--Polya program, we construct a compact integral realization
of the completed zeta function directly from its positive Theta kernel, while
separating this realization from the still-open problem of spectral reality.
Set $\Xi(z)=\xi(\tfrac12+iz)$, and let $W$ be the positive, even Theta density
whose Fourier transform is $\Xi$.  After normalizing
$p=W/\Xi(0)$ and writing $F(x)=\int_{-\infty}^x p(y)\dd y$, we introduce on
$L^2(\R)$ the centered Volterra operator
\[
 (\BTheta f)(x)=\frac{i}{\sqrt{p(x)}}\left(
 \int_{-\infty}^x\sqrt{p(y)}f(y)\dd y
 -F(x)\int_{\R}\sqrt{p(y)}f(y)\dd y\right).
\]
We prove that $\BTheta$ is Hilbert--Schmidt, obtain an explicit norm bound,
and establish the regularized determinant identity
\[
 \detTwo(I+z\BTheta)=\frac{\Xi(z)}{\Xi(0)},\qquad z\in\C.
\]
Within the natural class of scalar rank-one endpoint corrections, the CDF
term is uniquely forced by requiring differentiation to remove exactly the
$\sqrt p$ mode.  Thus the construction is canonical for the underlying
probability law rather than an ad hoc determinant reduction.
Consequently the nonzero spectrum of $\BTheta$ consists, with algebraic
multiplicity, of the reciprocal negatives of the zeros of $\Xi$.  We also
construct a closed first-order operator
$A=p^{-1/2}(-i\,d/dx)p^{1/2}$ and prove directly that
$\BTheta=A^\dagger$, its Moore--Penrose inverse.  The associated block
supercharge gives a one-dimensional Witten-type factorization, and $A$ is
Fredholm of index $-1$.  For an even probability density, $A$ and
$\BTheta$ are PT-symmetric and parity-odd.  The determinant formula is first
proved for an arbitrary density on a finite interval and is then extended to
an explicit class of whole-line densities by Hilbert--Schmidt convergence.
This yields, in particular, trace--cumulant identities for the associated
operators.  The construction is unconditional; it does not prove the Riemann
hypothesis.  Rather, the hypothesis becomes the precise additional assertion
that the nonzero spectrum of this concrete compact PT-symmetric operator is
real.
\end{abstract}

\maketitle

\section{Introduction}

Write
\begin{equation}\label{eq:xi-def}
 \xi(s)=\frac12s(s-1)\pi^{-s/2}\Gamma(s/2)\zeta(s),
 \qquad
 \Xi(z)=\xi\!\left(\frac12+iz\right).
\end{equation}
The Riemann hypothesis (RH) is equivalent to the assertion that every zero of
the even entire function $\Xi$ is real.  A traditional Hilbert--Polya program
therefore asks for an operator whose spectral points are the zeros of $\Xi$
and whose reality follows from self-adjointness.  Connes's noncommutative
geometric realization \cite{Connes1999} is a foundational example of a
broader spectral interpretation: the critical zeros occur as an absorption
spectrum for the idele-class action on the noncommutative space of Adele
classes, while the explicit formulas are interpreted through a trace formula.
Recent work of Hedenmalm
\cite{Hedenmalm2026} gives a Theta-based boundary-value formulation of the
zeros, including complex zeros, and isolates the search for a suitable inner
product.  The present paper takes a different but compatible step: it turns
the same positive Theta density into an explicit compact integral operator and
computes its regularized Fredholm determinant exactly.

Our main object is not obtained from a table of zeros.  It is constructed
directly from the classical Riemann Fourier kernel.  Let $p$ be the normalized
kernel, let $F$ be its cumulative distribution function, and define
\begin{equation}\label{eq:B-intro}
 (\BTheta f)(x)=\frac{i}{\sqrt{p(x)}}
 \left[\int_{-\infty}^x\sqrt{p(y)}f(y)\dd y
 -F(x)\int_{\R}\sqrt{p(y)}f(y)\dd y\right].
\end{equation}
The first term records the ordered Volterra history $y\leq x$.  The rank-one
second term is forced by the two endpoint conditions: it removes precisely the
outgoing scalar left by the first term.  More strongly,
Proposition~\ref{prop:CDF-uniqueness} shows that $F$ is the unique scalar correction for
which differentiation of the ordered primitive returns the orthogonal
projection off the $\sqrt p$ mode.  Thus CDF centering is not selected merely
to make a determinant calculation close: it is forced by a local
projection-and-endpoint principle.

It is useful to distinguish this construction from two more immediate
realizations of a characteristic function.  If $M_x$ denotes multiplication
by the coordinate on $L^2(\R)$, then every probability density has the unitary
matrix-coefficient representation
\begin{equation}\label{eq:matrix-coefficient}
 \varphi_p(z)=
 \left\langle e^{izM_x}\sqrt p,\sqrt p\right\rangle.
\end{equation}
Although canonical, \eqref{eq:matrix-coefficient} does not turn the zeros of
$\varphi_p$ into spectral points of a single compact operator.  At the other
extreme, a diagonal or canonical-product model assembled from a prescribed
zero set has the desired zeros by construction but uses precisely the data one
wants to explain.  The probability-CDF centering takes neither route: both the
Volterra part and its rank-one correction are determined by $p$ before any
zero is known, while the modified Fredholm determinant of the resulting one
operator recovers the complete centered characteristic function.  The CDF
term also retains the endpoint datum that a bare Volterra operator loses.

General semiseparable determinant theory supplies a broad reduction framework
for kernels of this type \cite{GesztesyMakarov2003}.  What is specific here is
the probability normalization that makes the scalar reduction close exactly
on a characteristic function, together with a whole-line Schatten criterion
that can be verified with explicit constants for the Riemann--Theta density.
Compared with the Theta boundary formulation in \cite{Hedenmalm2026}, the
present operator is bounded and compact, is constructed without selecting a
zero, and records the entire zero divisor and its algebraic multiplicities in
one regularized determinant.  These are complementary features rather than a
claim that one framework subsumes the other.

Two further distinctions locate the construction more precisely.  Connes's
spectral side is global and representation-theoretic: it is built from the
idele-class action on the Adele-class space, critical zeros are missing
spectral lines, and possible noncritical zeros enter as resonances
\cite{Connes1999}.  No Adele space, group representation, or arithmetic trace
formula is used here; the operator in \eqref{eq:B-intro} is a compact operator
on $L^2(\R)$ obtained directly from a positive probability density.  Suzuki
\cite{Suzuki2020}, in a different integral-operator approach, studies the
Hankel-type family $\mathsf K_{\zeta}^{\omega,\nu}[t]$ on
$L^2(-\infty,t)$.  Its kernel is defined through a shifted ratio of completed
zeta functions, and the quotient
\[
 \frac{\det(I+\mathsf K_{\zeta}^{\omega,\nu}[t])}
      {\det(I-\mathsf K_{\zeta}^{\omega,\nu}[t])}
\]
enters a de Branges Hamiltonian obtained from associated integral equations.
By contrast, the present construction uses neither de Branges spaces nor the
parameters $(\omega,\nu,t)$: a single fixed Hilbert--Schmidt operator
$\BTheta$ is determined by the classical Theta density before any zero is
chosen, and its one modified determinant is directly
$\Xi(z)/\Xi(0)$.  Thus the common Fredholm vocabulary conceals different
operator data, determinant identities, and spectral objectives.

The principal results are as follows.
\begin{enumerate}[label=(\roman*),leftmargin=2.3em]
 \item For a positive density on a finite interval, the CDF is the unique
 rank-one scalar centering compatible with the local projection principle,
 and the regularized determinant of the resulting operator is the
 characteristic function with its linear phase removed; see
 Proposition~\ref{prop:CDF-uniqueness} and
 Theorem~\ref{thm:finite-det}.
 \item A transparent whole-line admissibility condition gives
 Hilbert--Schmidt convergence of symmetric truncations and hence the same
 determinant formula on $\R$; see \cref{thm:whole-line-det}.
 \item The Riemann--Theta density satisfies these assumptions.  Explicit
 first-mode estimates give
 \begin{equation}\label{eq:intro-hs-bound}
  \|\BTheta\|_{\HS}^2
  \leq \frac{1+\Delta}{\pi d_*}<\infty,
 \end{equation}
 where $\Delta$ and $d_*$ are given in \cref{eq:Delta-def,eq:dstar-def}.
 \item The resulting determinant is exactly
 \begin{equation}\label{eq:intro-main-det}
  \detTwo(I+z\BTheta)=\frac{\Xi(z)}{\Xi(0)}.
 \end{equation}
 Thus the zero divisor, including algebraic multiplicities, is realized by a
 single compact operator.
 \item The same integral kernel is the Moore--Penrose inverse of a natural
 closed first-order Theta operator.  This gives exact inverse identities and
 an explicit positive lower bound without using the invalid implication
 ``injective implies bounded below.''  The first-order factor has a
 Witten-type supersymmetric completion and Fredholm index $-1$.
\end{enumerate}

Up to an overall phase and interchange of the two partner sectors, the
first-order factor is the one-dimensional Witten differential associated
with the logarithmic potential $\Phi=-\log p$ \cite{Witten1982}.  The
PT symmetry of the compact inverse belongs to the non-Hermitian framework
initiated in \cite{BenderBoettcher1998}, while the possible metric
intertwining equation is naturally read in the pseudo-Hermitian framework of
\cite{Mostafazadeh2002,Mostafazadeh2002III}.  These connections sharpen the
spectral question but do not supply spectral reality.

The logical boundary is important.  Equation \eqref{eq:intro-main-det} is an
operator realization of the full zero divisor, not a proof that the divisor
is real.  In fact,
\[
 \mathrm{RH}
 \quad\Longleftrightarrow\quad
 \Spec(\BTheta)\setminus\{0\}\subset\R.
\]
PT symmetry supplies the expected quartet symmetry, but PT-symmetric compact
operators need not have real spectrum.  The missing step is therefore a
specific spectral-reality theorem for the arithmetic kernel, not a domain
closure, determinant normalization, or numerical approximation.

The determinant and trace-ideal facts used below are standard; our conventions
follow Simon \cite{Simon2005}.  Accordingly, the finite-interval calculation
in \cref{thm:finite-det} is presented as a self-contained scalar
specialization, not as a claim that rank-one Volterra determinant reduction is
new.  The contribution here is the probability-CDF centering, the explicit
full-line Schatten closure for the Riemann density, its Moore--Penrose
interpretation, and the exact specialization to $\Xi$.  Standard properties
of $\xi$ and its Fourier kernel may be found in Titchmarsh's account
\cite{Titchmarsh1986}.

\section{The positive Riemann--Theta density}

For $x\geq0$, set
\begin{equation}\label{eq:W-def}
 W(x)=2\sum_{n=1}^{\infty}
 \left(2\pi^2n^4e^{9x/2}-3\pi n^2e^{5x/2}\right)
 e^{-\pi n^2e^{2x}},
 \qquad W(-x)=W(x).
\end{equation}
The Jacobi transformation law shows that this even extension is smooth.  It
is the logarithmic form of the classical Theta density appearing in the
Mellin representation of $\xi$.  For completeness, a derivation from the
classical Jacobi inversion and Mellin formulas, independent of any recent
Theta-based spectral construction, is given in
\cref{sec:appendix-kernel}; compare \cite{Titchmarsh1986}.

\begin{proposition}[Riemann Fourier representation]\label{prop:fourier-xi}
The function $W$ is strictly positive, even, smooth, and has exponential
moments of every order.  Moreover,
\begin{equation}\label{eq:fourier-xi}
 \Xi(z)=\int_{\R}W(x)e^{izx}\dd x,
 \qquad z\in\C.
\end{equation}
In particular,
\begin{equation}\label{eq:c0-def}
 c_0:=\int_{\R}W(x)\dd x=\Xi(0)>0.
\end{equation}
\end{proposition}

\begin{proof}
The complete calculation is recorded in \cref{sec:appendix-kernel}.  In
brief, if $\vartheta(t)=\sum_{n\in\mathbb Z}e^{-\pi n^2t}$ and
\[
 \omega(t)=\frac12\bigl(2t^2\vartheta''(t)+3t\vartheta'(t)\bigr),
\]
then Jacobi inversion gives
$\omega(1/t)=t^{1/2}\omega(t)$, and
$W(x)=2e^{x/2}\omega(e^{2x})$.  This identity proves smoothness and evenness
without gluing two half-line formulas.  The series is positive for $x\geq0$
because $2\pi n^2e^{2x}-3>0$, hence it is positive everywhere by evenness,
and it has double-exponential decay.  Two integrations by parts in the
classical Mellin formula give
\[
 \xi(s)=\int_0^\infty\omega(t)t^{s/2}\frac{\dd t}{t}.
\]
The substitution $t=e^{2x}$ at $s=\tfrac12+iz$ is exactly
\eqref{eq:fourier-xi}.  Double-exponential decay gives absolute convergence
locally uniformly in $z\in\C$.
\end{proof}

Henceforth put
\begin{equation}\label{eq:pF-def}
 p(x)=\frac{W(x)}{c_0},
 \qquad
 F(x)=\int_{-\infty}^xp(y)\dd y.
\end{equation}
Then $p$ is an even, strictly positive probability density, $F(-x)=1-F(x)$,
and
\begin{equation}\label{eq:char-xi}
 \varphi_p(z):=\int_{\R}p(x)e^{izx}\dd x
 =\frac{\Xi(z)}{\Xi(0)}.
\end{equation}

\section{The finite-interval determinant identity}

We first isolate the operator-theoretic mechanism from the arithmetic
specialization.  Throughout, Hilbert-space inner products are linear in the
first variable.  For vectors $u,v$, the notation $u\otimes v$ means
$(u\otimes v)f=\langle f,v\rangle u$.

Let $-\infty<a<b<\infty$, and let $p$ be a probability density on $[a,b]$
that is bounded above and bounded away from zero.  Define
\begin{equation}\label{eq:finite-Fmu}
 F(x)=\int_a^xp(y)\dd y,
 \qquad
 \mu=\int_a^b xp(x)\dd x.
\end{equation}

\begin{definition}[Centered Volterra operator]\label{def:centered-volterra}
The centered Volterra operator associated with $p$ is
\begin{equation}\label{eq:finite-B}
 (\Vp f)(x)=\frac{i}{\sqrt{p(x)}}\left[
 \int_a^x\sqrt{p(y)}f(y)\dd y
 -F(x)\int_a^b\sqrt{p(y)}f(y)\dd y\right].
\end{equation}
\end{definition}

The particular correction $F$ in \eqref{eq:finite-B} is canonical within a
natural class.  Put $P_p=I-\sqrt p\otimes\sqrt p$.  For an absolutely
continuous function $h$ with $h(a)=0$, define
\begin{equation}\label{eq:Bh-def}
 (\mathcal B_{p,h}f)(x)=\frac{i}{\sqrt{p(x)}}\left[
 \int_a^x\sqrt{p(y)}f(y)\dd y
 -h(x)\int_a^b\sqrt{p(y)}f(y)\dd y\right].
\end{equation}

\begin{proposition}[Uniqueness of CDF centering]\label{prop:CDF-uniqueness}
The local projection identity
\begin{equation}\label{eq:CDF-local-projection}
 \frac{d}{dx}\left(\frac{\sqrt p\,\mathcal B_{p,h}f}{i}\right)
 =\sqrt p\,P_pf
 \qquad\text{for every }f\in L^2([a,b])
\end{equation}
holds in the weak sense if and only if $h=F$.  For this unique choice, the
bracket in \eqref{eq:Bh-def} vanishes at both $a$ and $b$ for every $f$.
\end{proposition}

\begin{proof}
Writing $m_f=\langle f,\sqrt p\rangle$, the two sides of
\eqref{eq:CDF-local-projection} are respectively
\[
 \sqrt p\,f-h'm_f
 \qquad\text{and}\qquad
 \sqrt p\,f-pm_f.
\]
Taking $f=\sqrt p$ shows that the identity holds for every $f$ only if
$h'=p$ almost everywhere.  Since $h(a)=0$, this is equivalent to $h=F$;
the converse is immediate.  Finally, $F(a)=0$ and $F(b)=1$, so the bracket
vanishes at both endpoints.
\end{proof}

The operator $\Vp$ in Definition~\ref{def:centered-volterra} is
Hilbert--Schmidt.
Conjugation by multiplication with $\sqrt p$ turns it into
\begin{equation}\label{eq:C-similar}
 C=i(\mathcal V-F\otimes\1),
 \qquad
 (\mathcal V f)(x)=\int_a^xf(y)\dd y,
\end{equation}
on $L^2([a,b])$.  The multiplication operators are bounded and boundedly
invertible, so the regularized determinant is unchanged by this similarity.

For $K\in\HS$, we use
\begin{equation}\label{eq:det2-def}
 \detTwo(I+K)=\det\bigl((I+K)e^{-K}\bigr).
\end{equation}
It is an entire function of scalar multiples of $K$ and is continuous in the
Hilbert--Schmidt norm \cite[Chapter 9]{Simon2005}.  We use the standard
Hilbert--Schmidt product formula in the form recorded, for example, in
\cite{BritzEtAl2021}.

\begin{theorem}[Finite centered-Volterra determinant]\label{thm:finite-det}
For every $z\in\C$,
\begin{equation}\label{eq:finite-det}
 \boxed{
 \detTwo(I+z\Vp)
 =e^{-i\mu z}\int_a^bp(x)e^{izx}\dd x.}
\end{equation}
\end{theorem}

\begin{proof}
The Volterra operator $\mathcal V$ is Hilbert--Schmidt and quasinilpotent;
hence
\begin{equation}\label{eq:volterra-det2-one}
 \detTwo(I+iz\mathcal V)=1.
\end{equation}
Put
\[
 Q=(I+iz\mathcal V)^{-1}F.
\]
Since $\mathcal V$ commutes with its resolvent, the factorization
\begin{equation}\label{eq:finite-factor}
 I+zC=(I+iz\mathcal V)(I-izQ\otimes\1)
\end{equation}
holds.  The multiplicative formula for $\detTwo$ gives
\begin{equation}\label{eq:det-step1}
\begin{aligned}
 \detTwo(I+zC)
 &=\biggl(1-iz\int_a^bQ(x)\dd x\biggr)
 \exp\biggl(
 iz\int_a^bQ(x)\dd x\\
 &\hspace{8em}
 -z^2\int_a^b(\mathcal VQ)(x)\dd x\biggr).
\end{aligned}
\end{equation}
The resolvent identity $(I+iz\mathcal V)Q=F$ shows that the exponent in
\eqref{eq:det-step1} is $iz\int_a^bF$.

For the scalar factor, set
\[
 q=(I+iz\mathcal V)^{-1}p.
\]
Because $F=\mathcal Vp$, one has $Q=\mathcal Vq$.  Therefore
\begin{equation}\label{eq:bracket-Qb}
 1-iz\int_a^bQ=\int_a^bq=Q(b).
\end{equation}
The function $Q$ satisfies the initial-value problem
\begin{equation}\label{eq:Q-ode}
 Q'(x)+izQ(x)=p(x),
 \qquad Q(a)=0,
\end{equation}
and consequently
\begin{equation}\label{eq:Qb-fourier}
 Q(b)=e^{-izb}\int_a^bp(x)e^{izx}\dd x.
\end{equation}
Finally, Fubini's theorem gives
\begin{equation}\label{eq:intF}
 \int_a^bF(x)\dd x
 =\int_a^b(b-y)p(y)\dd y=b-\mu.
\end{equation}
Combining \eqref{eq:det-step1}--\eqref{eq:intF} proves
\eqref{eq:finite-det}.  Similarity in \eqref{eq:C-similar} returns the result
to $\Vp$.
\end{proof}

\begin{corollary}[Finite trace--cumulant identity]\label{cor:finite-cumulant}
Let $\kappa_n(p)$ be the cumulants of $p$.  For every $n\geq2$,
\begin{equation}\label{eq:finite-cumulant}
 \tr(\Vp^n)
 =\frac{(-1)^{n+1}i^n}{(n-1)!}\,\kappa_n(p).
\end{equation}
\end{corollary}

\begin{proof}
Near $z=0$, compare
\[
 \log\detTwo(I+z\Vp)
 =\sum_{n=2}^{\infty}
 \frac{(-1)^{n+1}}{n}z^n\tr(\Vp^n)
\]
with the cumulant expansion of
$e^{-i\mu z}\int p(x)e^{izx}\dd x$.
\end{proof}

\section{Whole-line admissibility and closure}

The finite-interval formula does not pass automatically to $\R$: the
centered Volterra kernel may fail to be Hilbert--Schmidt, even for familiar
density functions.  We now state a directly checkable class for which the
passage is valid.

Let $p$ be a strictly positive, continuous, even probability density on
$\R$, and let $F(x)=\int_{-\infty}^xp$.  Define formally
\begin{equation}\label{eq:whole-B}
 (\Vp f)(x)=\frac{i}{\sqrt{p(x)}}\left[
 \int_{-\infty}^x\sqrt{p(y)}f(y)\dd y
 -F(x)\int_{\R}\sqrt{p(y)}f(y)\dd y\right].
\end{equation}

\begin{definition}[Volterra-admissible density]\label{def:admissible}
An even density $p$ is called \emph{Volterra-admissible} if:
\begin{enumerate}[label=(A\arabic*),leftmargin=3.1em]
 \item\label{item:A1} $p(x)>0$ for all $x$ and
 $\int_{\R}p(x)e^{M|x|}\dd x<\infty$ for every $M>0$;
 \item\label{item:A2}
 \begin{equation}\label{eq:admissible-HS}
  I_p:=\int_{\R}\frac{F(x)(1-F(x))}{p(x)}\dd x<\infty;
 \end{equation}
 \item\label{item:A3} if
 $q_R=\int_R^\infty p(x)\dd x$, then
 \begin{equation}\label{eq:admissible-defect}
  \frac{q_R^2}{1-2q_R}
  \int_{-R}^R\frac{(2F(x)-1)^2}{p(x)}\dd x
  \longrightarrow0.
 \end{equation}
\end{enumerate}
\end{definition}

The terminology is local to this paper.  Condition \eqref{eq:admissible-HS}
is exactly the Hilbert--Schmidt requirement; condition
\eqref{eq:admissible-defect} controls the normalization error created by
symmetric truncation.

\begin{proposition}[Exact Hilbert--Schmidt norm]\label{prop:HS-identity}
If \eqref{eq:admissible-HS} holds, then \eqref{eq:whole-B} defines an operator
in $\HS(L^2(\R))$ with kernel
\begin{equation}\label{eq:whole-kernel}
 b_p(x,y)=i\sqrt{\frac{p(y)}{p(x)}}
 \bigl(\1_{\{y\leq x\}}-F(x)\bigr)
\end{equation}
and
\begin{equation}\label{eq:HS-exact}
 \boxed{\|\Vp\|_{\HS}^2
 =\int_{\R}\frac{F(x)(1-F(x))}{p(x)}\dd x.}
\end{equation}
\end{proposition}

\begin{proof}
Tonelli's theorem applies to $|b_p|^2$.  For fixed $x$,
\begin{align*}
 \int_{\R}|b_p(x,y)|^2\dd y
 &=\frac{1}{p(x)}\left[(1-F(x))^2
 \int_{-\infty}^xp(y)\dd y
 +F(x)^2\int_x^\infty p(y)\dd y\right]\\
 &=\frac{F(x)(1-F(x))}{p(x)}.
\end{align*}
Integrating the preceding identity over all $x\in\R$ gives
\eqref{eq:HS-exact} and proves the claim.
\end{proof}

For $R>0$, define
\begin{equation}\label{eq:pR-def}
 p_R(x)=\frac{p(x)\1_{[-R,R]}(x)}{1-2q_R},
 \qquad
 F_R(x)=\int_{-R}^xp_R(y)\dd y,
\end{equation}
and let $\Vp^{(R)}$ be the finite-interval operator, extended by zero to
$L^2(\R)$.

\begin{proposition}[Hilbert--Schmidt truncation]\label{prop:HS-truncation}
If $p$ is Volterra-admissible, then
\begin{equation}\label{eq:HS-truncation}
 \|\Vp^{(R)}-\Vp\|_{\HS}\longrightarrow0.
\end{equation}
\end{proposition}

\begin{proof}
On $[-R,R]$ one has the exact identity
\begin{equation}\label{eq:FR-F}
 F_R(x)-F(x)=\frac{q_R(2F(x)-1)}{1-2q_R}.
\end{equation}
The normalization in $p_R$ cancels from the ratio
$\sqrt{p_R(y)/p_R(x)}$.  Therefore the squared kernel error on
$[-R,R]^2$ is exactly
\begin{equation}\label{eq:interior-error}
 E_R^{\mathrm{in}}
 =\frac{q_R^2}{1-2q_R}
 \int_{-R}^R\frac{(2F(x)-1)^2}{p(x)}\dd x,
\end{equation}
which tends to zero by \eqref{eq:admissible-defect}.  On the complement of
the square, the truncated kernel vanishes.  The remaining error is the
integral of $|b_p|^2$ over a decreasing family of sets with empty intersection;
it tends to zero by \cref{prop:HS-identity} and continuity from above.
\end{proof}

\begin{theorem}[Whole-line determinant]\label{thm:whole-line-det}
If $p$ is Volterra-admissible, then
\begin{equation}\label{eq:whole-line-det}
 \boxed{
 \detTwo(I+z\Vp)=\int_{\R}p(x)e^{izx}\dd x,
 \qquad z\in\C.}
\end{equation}
The convergence of the finite determinants to both sides is locally uniform
on $\C$.
\end{theorem}

\begin{proof}
Each $p_R$ is even, so its mean is zero.  By \cref{thm:finite-det},
\[
 \detTwo(I+z\Vp^{(R)})
 =\int_{-R}^Rp_R(x)e^{izx}\dd x.
\]
The left side converges locally uniformly by
\cref{prop:HS-truncation} and continuity estimates for $\detTwo$ on $\HS$.
For $|\operatorname{Im}z|\leq M$, the right side converges uniformly by
condition \ref{item:A1}, normalization of $p_R$, and dominated convergence.
The two limits agree.
\end{proof}

\begin{remark}[Nonzero mean]\label{rem:nonzero-mean}
For a non-even density of mean $\mu$, the same argument, with two-sided
truncation defects written separately, gives
\[
 \detTwo(I+z\Vp)=e^{-i\mu z}\varphi_p(z).
\]
Thus the regularization removes exactly the first cumulant.
\end{remark}

\begin{example}\label{ex:supergaussian}
The densities $p_r(x)=c_r e^{-|x|^r}$ are Volterra-admissible for $r>2$.
Indeed, their tail-to-density ratio is $O(x^{1-r})$, which is integrable
precisely beyond the Gaussian exponent.  At $r=2$, the Gaussian density fails
\eqref{eq:admissible-HS}: Mills' ratio gives
$F(x)(1-F(x))/p(x)\sim(2x)^{-1}$ as $x\to+\infty$.  Hence whole-line
Hilbert--Schmidt closure is a genuine restriction, not a formal consequence
of rapid decay.
\end{example}

\begin{remark}[The reciprocal log-slope test]\label{rem:log-slope}
The Gaussian obstruction is not a lack of moments: the Gaussian has
exponential moments of every order.  The obstruction comes from the
$p^{-1/2}$ normalization in the kernel.  On the positive tail,
$F(x)\to1$, so the Hilbert--Schmidt integrand is asymptotic to
$(1-F(x))/p(x)$.  For a smooth, eventually log-convex tail
$p(x)=c\exp(-\Phi(x))$ satisfying $\Phi'(x)\to\infty$ and
$\Phi''(x)/(\Phi'(x))^2\to0$, endpoint Laplace estimates give
\begin{equation}\label{eq:log-slope-diagnostic}
 \frac{1-F(x)}{p(x)}\sim\frac1{\Phi'(x)}.
\end{equation}
Thus admissibility tests the integrability of the reciprocal logarithmic
slope, not merely rapid decay.  For $\Phi(x)=x^2$, the right side is
$(2x)^{-1}$ and diverges logarithmically.  By contrast,
\cref{eq:w1-def,eq:W-w1-comparison} show that the Riemann--Theta density has
\begin{equation}\label{eq:theta-log-slope}
 -\log p(x)=\pi e^{2x}-\frac92x+O(1),
 \qquad x\to+\infty,
\end{equation}
so its reciprocal log-slope is of order $e^{-2x}$.  The rigorous version is
the tail ratio bound \eqref{eq:tail-density-ratio}.  This double-exponential geometry
is exactly what separates the Theta density from the Gaussian threshold.
\end{remark}

\section{Explicit Theta-tail estimates}

We now verify Definition~\ref{def:admissible} for the density in
\cref{eq:pF-def} and
derive a fixed-constant norm bound.

For $x\geq0$, put
\begin{equation}\label{eq:q-def}
 q=\pi e^{2x}\geq\pi
\end{equation}
and denote the $n=1$ summand of \eqref{eq:W-def} by
\begin{equation}\label{eq:w1-def}
 w_1(x)=2\pi^{-1/4}q^{5/4}(2q-3)e^{-q}.
\end{equation}
For $n\geq2$, the ratio of the $n$th summand to $w_1$ is
\begin{equation}\label{eq:wn-ratio}
 \frac{w_n(x)}{w_1(x)}
 =n^2\frac{2qn^2-3}{2q-3}e^{-q(n^2-1)}
 \leq2n^4e^{-q(n^2-1)}.
\end{equation}
Since $n^2-1\geq3(n-1)$, define
\begin{equation}\label{eq:Delta-def}
 r=e^{-3\pi},
 \qquad
 \Delta=2\left[
 \frac{1+11r+11r^2+r^3}{(1-r)^5}-1\right].
\end{equation}
Summing \eqref{eq:wn-ratio} gives
\begin{equation}\label{eq:W-w1-comparison}
 \boxed{w_1(x)\leq W(x)\leq(1+\Delta)w_1(x),
 \qquad x\geq0.}
\end{equation}

We use the elementary incomplete-Gamma estimate
\begin{equation}\label{eq:incomplete-gamma}
 \int_q^\infty t^ae^{-t}\dd t
 \leq\frac{q^ae^{-q}}{1-a/q},
 \qquad q>a\geq0.
\end{equation}
It follows from
$t^a\leq q^a\exp(a(t-q)/q)$ for $t\geq q$.

Let $T(x)=\int_x^\infty W(y)\dd y$.  The substitution
$t=\pi e^{2y}$ yields
\begin{equation}\label{eq:w1-tail-change}
 \int_x^\infty w_1(y)\dd y
 =\pi^{-1/4}\int_q^\infty
 t^{1/4}(2t-3)e^{-t}\dd t.
\end{equation}
Combining \eqref{eq:W-w1-comparison}--\eqref{eq:w1-tail-change} gives
\begin{equation}\label{eq:tail-density-ratio}
 \frac{T(x)}{W(x)}
 \leq\frac{1+\Delta}{(2q-3)(1-5/(4q))}
 \leq\frac{1+\Delta}{\pi d_*}e^{-2x},
\end{equation}
where
\begin{equation}\label{eq:dstar-def}
 d_*=\left(2-\frac3\pi\right)
 \left(1-\frac5{4\pi}\right)>0.
\end{equation}

\begin{theorem}[Theta admissibility and norm bound]\label{thm:theta-admissible}
The density $p=W/c_0$ is Volterra-admissible.  The corresponding operator
$\BTheta:=\Vp$ satisfies
\begin{equation}\label{eq:theta-HS-bound}
 \boxed{
 \|\BTheta\|_{\HS}^2
 \leq\frac{1+\Delta}{\pi d_*}.}
\end{equation}
\end{theorem}

\begin{proof}
For $x\geq0$, $c_0(1-F(x))=T(x)$ and $F(x)\leq1$.  By evenness and
\eqref{eq:tail-density-ratio},
\begin{align}
 \|\BTheta\|_{\HS}^2
 &=c_0\int_{\R}\frac{F(x)(1-F(x))}{W(x)}\dd x\notag\\
 &\leq2\int_0^\infty\frac{T(x)}{W(x)}\dd x
 \leq\frac{1+\Delta}{\pi d_*}.
 \label{eq:theta-HS-proof}
\end{align}
This proves condition \ref{item:A2}.

It remains to control symmetric truncation.  Put
\begin{equation}\label{eq:QR-Cq}
 Q_R=\pi e^{2R},
 \qquad
 C_q=\frac{2(1+\Delta)\pi^{-1/4}}
 {c_0(1-5/(4\pi))}.
\end{equation}
Equations \eqref{eq:W-w1-comparison} and
\eqref{eq:incomplete-gamma} imply
\begin{equation}\label{eq:qR-bound}
 q_R\leq C_qQ_R^{5/4}e^{-Q_R}.
\end{equation}
The lower bound $W\geq w_1$ gives
\begin{equation}\label{eq:reciprocal-bound}
 \int_{-R}^R\frac{\dd x}{p(x)}
 \leq\frac{c_0}{2\pi^2(2\pi-3)}e^{Q_R}.
\end{equation}
For all sufficiently large $R$, $q_R\leq1/4$; using
$(2F-1)^2\leq1$, the expression in
\eqref{eq:admissible-defect} is at most
\begin{equation}\label{eq:explicit-interior-bound}
 \frac{C_q^2c_0}{\pi^2(2\pi-3)}
 Q_R^{5/2}e^{-Q_R},
\end{equation}
which tends to zero.  This proves condition \ref{item:A3}.

Finally, inserting $e^{M|x|}$ in the same calculation, with
$a_M=5/4+M/2$ and $Q_R>a_M$, gives
\begin{equation}\label{eq:exp-tail-bound}
 \int_{|x|>R}p(x)e^{M|x|}\dd x
 \leq
 \frac{4(1+\Delta)\pi^{-1/4-M/2}}
 {c_0(1-a_M/Q_R)}Q_R^{a_M}e^{-Q_R}.
\end{equation}
Thus condition \ref{item:A1} also holds.
\end{proof}

\begin{corollary}[Fredholm realization of $\Xi$]\label{cor:main-det}
For every $z\in\C$,
\begin{equation}\label{eq:main-det}
 \boxed{
 \detTwo(I+z\BTheta)=\frac{\Xi(z)}{\Xi(0)}.}
\end{equation}
\end{corollary}

\begin{proof}
Combine \cref{eq:char-xi,thm:whole-line-det,thm:theta-admissible}.
\end{proof}

\section{A first-order operator and its Moore--Penrose inverse}

The integral operator also has a differential interpretation.  Let
\begin{equation}\label{eq:psi-P0}
 \psi=\sqrt p,
 \qquad
 P_0=I-\psi\otimes\psi.
\end{equation}
Since $p$ is normalized, $\|\psi\|_2=1$.  Put
\begin{equation}\label{eq:D-def}
 D=-i\frac{d}{dx}
\end{equation}
and define the maximal realization
\begin{align}
 \Dom(A)=\{g\in L^2(\R):{}&u=\sqrt p\,g\in H^1(\R),
 \ p^{-1/2}Du\in L^2(\R)\},\notag\\
 Ag={}&p^{-1/2}D(\sqrt p\,g).
 \label{eq:A-domain}
\end{align}
The space $C_c^\infty(\R)$ is contained in $\Dom(A)$, so $A$ is densely
defined.

\begin{proposition}\label{prop:A-closed}
The operator $A$ is closed.  If $\Phi=-\log p$ and
$V=\tfrac12\Phi'$, then on $C_c^\infty(\R)$,
\begin{equation}\label{eq:A-DV}
 A=D+iV.
\end{equation}
\end{proposition}

\begin{proof}
Suppose $g_n\to g$ and $Ag_n\to h$ in $L^2$.  Since $\sqrt p$ is bounded,
$\sqrt p\,g_n\to\sqrt p\,g$ and
$D(\sqrt p\,g_n)=\sqrt p\,Ag_n\to\sqrt p\,h$ in $L^2$.  Closedness of
$D$ on $H^1(\R)$ gives $g\in\Dom(A)$ and $Ag=h$.  Formula
\eqref{eq:A-DV} is a direct product-rule calculation.
\end{proof}

\begin{theorem}[Explicit Moore--Penrose inverse]\label{thm:MP-inverse}
The following identities hold:
\begin{align}
 A\BTheta&=P_0 &&\text{on }L^2(\R),
 \label{eq:AB-P}\\
 \BTheta A&=I &&\text{on }\Dom(A),
 \label{eq:BA-I}\\
 \BTheta\psi&=0.\label{eq:Bpsi-zero}
\end{align}
Consequently,
\begin{equation}\label{eq:B-MP}
 \boxed{\BTheta=A^\dagger,}
\end{equation}
$A$ is injective, $\Ran(A)=\psi^\perp$ is closed, and
\begin{equation}\label{eq:A-lower}
 \boxed{
 \|Ag\|_2\geq
 \sqrt{\frac{\pi d_*}{1+\Delta}}\,\|g\|_2,
 \qquad g\in\Dom(A).}
\end{equation}
\end{theorem}

\begin{proof}
For $f\in L^2(\R)$, set
\begin{equation}\label{eq:qf-def}
 m=\int_{\R}\sqrt{p(y)}f(y)\dd y,
 \qquad
 q_f(x)=\int_{-\infty}^x\sqrt{p(y)}f(y)\dd y-F(x)m.
\end{equation}
Then $\sqrt p\,\BTheta f=iq_f$ and
\begin{equation}\label{eq:qf-prime}
 q_f'(x)=\sqrt{p(x)}(P_0f)(x).
\end{equation}
The Hilbert--Schmidt bound shows that $\BTheta f\in L^2$; hence
$q_f\in L^2$.  Equation \eqref{eq:qf-prime} puts $q_f$ in $H^1$ and gives
$A\BTheta f=P_0f$.

Conversely, let $g\in\Dom(A)$ and put $u=\sqrt p\,g$.  Every $H^1(\R)$
function has a continuous representative tending to zero at both endpoints.
If $f=Ag$, then
\begin{equation}\label{eq:Du-f}
 \sqrt p\,f=Du=-iu'.
\end{equation}
Moreover, $\sqrt p\,f\in L^1$ by Cauchy--Schwarz.  Therefore
\begin{equation}\label{eq:integrate-Du}
 \int_{-\infty}^x\sqrt p\,f=-iu(x),
 \qquad
 \int_{\R}\sqrt p\,f=0.
\end{equation}
Substitution in \eqref{eq:B-intro} gives $\BTheta Ag=g$.  Taking
$f=\psi$ in \eqref{eq:qf-def} gives $q_f=0$, hence
$\BTheta\psi=0$.

The three identities imply that $A$ is injective and has range exactly
$\psi^\perp$, while $\BTheta$ vanishes on the orthogonal complement of the
range.  These are the Moore--Penrose equations, proving
\eqref{eq:B-MP}.  Finally,
\[
 \|g\|_2=\|\BTheta Ag\|_2
 \leq\|\BTheta\|\,\|Ag\|_2
 \leq\|\BTheta\|_{\HS}\,\|Ag\|_2,
\]
and \eqref{eq:A-lower} follows from \eqref{eq:theta-HS-bound}.
\end{proof}

\begin{corollary}[Witten-type factorization and index]
\label{cor:Witten-factorization}
Let $q=iA$ and define, on the natural block domain,
\begin{equation}\label{eq:supercharge}
 \mathcal Q=
 \begin{pmatrix}0&q^*\\ q&0\end{pmatrix},
 \qquad
 \Gamma=
 \begin{pmatrix}I&0\\0&-I\end{pmatrix}.
\end{equation}
Then $\mathcal Q$ is self-adjoint, $\Gamma\mathcal Q=-\mathcal Q\Gamma$, and
\begin{equation}\label{eq:Witten-Hamiltonian}
 \mathcal Q^2=
 \begin{pmatrix}A^*A&0\\0&AA^*\end{pmatrix}.
\end{equation}
On $C_c^\infty(\R)$,
\begin{align}
 q&=\frac{d}{dx}-V,&q^*&=-\frac{d}{dx}-V,\notag\\
 q^*q&=-\frac{d^2}{dx^2}+V^2+V',&
 qq^*&=-\frac{d^2}{dx^2}+V^2-V'.
 \label{eq:Witten-partners}
\end{align}
Moreover,
\begin{equation}\label{eq:A-index}
 \Ker(q)=\{0\},\qquad
 \Ker(q^*)=\operatorname{span}\{\psi\},\qquad
 \operatorname{ind}A=\operatorname{ind}q=-1.
\end{equation}
\end{corollary}

\begin{proof}
The closed densely defined operator $q=iA$ produces the standard self-adjoint
block supercharge in \eqref{eq:supercharge}; its square and grading relation
are immediate.  Formula \eqref{eq:Witten-partners} follows from
\eqref{eq:A-DV}.  By \cref{thm:MP-inverse}, $\Ker(A)=\{0\}$ and
$\Ran(A)=\psi^\perp$ is closed.  Hence
$\Ker(A^*)=\Ran(A)^\perp=\operatorname{span}\{\psi\}$, which proves
\eqref{eq:A-index}.
\end{proof}

Up to an overall phase and the conventional interchange of partner sectors,
\cref{eq:supercharge,eq:Witten-partners} are the one-dimensional Witten
supersymmetric factorization with Witten function $-\Phi/2$ (equivalently,
superpotential derivative $-V$)
\cite{Witten1982}.  In this language $\psi=\sqrt p$ is the unique zero mode
in one partner sector, while $\BTheta=A^\dagger$ is the bounded
Moore--Penrose inverse of the first-order factor on $\psi^\perp$.  It is not a
resolvent of either self-adjoint partner Hamiltonian, and supersymmetry alone
does not imply that the nonnormal compact operator $\BTheta$ has real
spectrum.

\begin{remark}\label{rem:no-kernel-shortcut}
The order of proof matters.  Injectivity of a closed operator does not imply a
positive lower bound.  Here the valid chain is
\begin{align*}
 \text{explicit tail estimate}
 &\Longrightarrow \BTheta\in\HS\\
 &\Longrightarrow (A\BTheta=P_0,\ \BTheta A=I)\\
 &\Longrightarrow |A|\geq\|\BTheta\|^{-1}I.
\end{align*}
\end{remark}

\section{PT symmetry and the zero spectrum}

Let parity $\mathsf P$ and conjugation $\mathsf T$ act by
\begin{equation}\label{eq:PT-def}
 (\mathsf Pf)(x)=f(-x),
 \qquad
 (\mathsf Tf)(x)=\overline{f(x)}.
\end{equation}
Because $p$ is even, $V$ in \eqref{eq:A-DV} is real and odd.

\begin{proposition}[Symmetries]\label{prop:PT}
On their natural domains,
\begin{align}
 \mathsf P A\mathsf P&=-A,
 &\mathsf T A\mathsf T&=-A,
 &(\mathsf P\mathsf T)A(\mathsf P\mathsf T)&=A,
 \label{eq:A-PT}\\
 \mathsf P\BTheta\mathsf P&=-\BTheta,
 &(\mathsf P\mathsf T)\BTheta(\mathsf P\mathsf T)&=\BTheta.
 \label{eq:B-PT}
\end{align}
Thus $\BTheta$ is parity-odd and PT-symmetric.
\end{proposition}

\begin{proof}
The relations for $A=D+iV$ follow from
$\mathsf PDP=-D$, $\mathsf T D\mathsf T=-D$, and the oddness of $V$.
The relations for $\BTheta=A^\dagger$ follow by applying the same unitary and
antiunitary symmetries to the Moore--Penrose equations.  They can also be read
directly from \eqref{eq:whole-kernel} and $F(-x)=1-F(x)$.
\end{proof}

This is the standard antilinear symmetry setting of PT-symmetric quantum
theory \cite{BenderBoettcher1998}.  The distinction between symmetry and
spectral reality is essential: PT symmetry permits real eigenvalues but does
not force them.  In the pseudo-Hermitian formulation
\cite{Mostafazadeh2002,Mostafazadeh2002III}, one instead seeks a metric
operator $G$ satisfying $G\BTheta=\BTheta^*G$.  If $G$ is bounded, positive,
and boundedly invertible, then $G^{1/2}\BTheta G^{-1/2}$ is self-adjoint, so
the spectrum is real.  Conversely, abstract metric reconstructions require
additional diagonalizability and basis hypotheses; they cannot be imported
here from PT symmetry alone.

We now state the exact spectral consequence of the determinant formula.

\begin{theorem}[Reciprocal zero spectrum]\label{thm:zero-spectrum}
The operator $\BTheta$ is compact and
\begin{equation}\label{eq:zero-spectrum}
 \boxed{
 \Spec(\BTheta)\setminus\{0\}
 =\left\{-\frac1\alpha:\ \Xi(\alpha)=0\right\}.}
\end{equation}
Algebraic multiplicities agree: the order of a zero $\alpha$ of $\Xi$ equals
the algebraic multiplicity of $-1/\alpha$ as an eigenvalue of $\BTheta$.
Moreover,
\begin{equation}\label{eq:RH-equivalence}
 \boxed{
 \mathrm{RH}
 \quad\Longleftrightarrow\quad
 \Spec(\BTheta)\setminus\{0\}\subset\R.}
\end{equation}
\end{theorem}

\begin{proof}
Hilbert--Schmidt operators are compact.  For a compact operator $K$, the zeros
of $\detTwo(I+zK)$ are precisely the points $z=-1/\lambda$, where
$\lambda\neq0$ is an eigenvalue of $K$; the zero order is the algebraic
multiplicity \cite[Chapter 9]{Simon2005}.  Apply this fact to
\cref{cor:main-det}.  Finally, a nontrivial zeta zero
$\rho=\beta+i\gamma$ corresponds to
\begin{equation}\label{eq:alpha-rho}
 \alpha=\gamma-i\left(\beta-\frac12\right),
 \qquad
 \rho=\frac12+i\alpha.
\end{equation}
Thus $\beta=1/2$ if and only if $\alpha$, and hence $-1/\alpha$, is real.
\end{proof}

The symmetry relations imply that each nonzero spectral point occurs in the
quartet
\begin{equation}\label{eq:spectral-quartet}
 \lambda,\quad-\lambda,\quad\overline\lambda,
 \quad-\overline\lambda,
\end{equation}
with collapses on the axes.  This is the operator form of the evenness,
reality, and functional-equation symmetries of $\Xi$.

The connection with the Theta boundary state can also be written directly.
For $\alpha\in\C$, define
\begin{equation}\label{eq:u-alpha}
 u_\alpha(x)=e^{-i\alpha x}
 \int_{-\infty}^xe^{i\alpha y}W(y)\dd y,
 \qquad
 g_\alpha=\frac{u_\alpha}{\sqrt W}.
\end{equation}

\begin{lemma}[Boundary-state domain lemma]\label{lem:boundary-domain}
If $\Xi(\alpha)=0$, then $g_\alpha\in\Dom(A)$.  More precisely, there are
constants $C_\alpha,R_\alpha>0$ such that
\begin{equation}\label{eq:u-alpha-tail}
 |u_\alpha(x)|\leq C_\alpha e^{-2|x|}W(x),
 \qquad |x|\geq R_\alpha.
\end{equation}
\end{lemma}

\begin{proof}
The zero condition and \cref{eq:fourier-xi} give the second representation
\[
 u_\alpha(x)=-e^{-i\alpha x}
 \int_x^\infty e^{i\alpha y}W(y)\dd y.
\]
Use this formula as $x\to+\infty$ and the original formula as
$x\to-\infty$.  After taking absolute values, the only additional factor
relative to \eqref{eq:w1-tail-change} is
$e^{|\Im\alpha|\,|y-x|}$.  Under the substitution
$t=\pi e^{2|y|}$ it changes the power of $t$ by
$|\Im\alpha|/2$.  The incomplete-Gamma estimate
\eqref{eq:incomplete-gamma}, together with
\eqref{eq:W-w1-comparison}, therefore gives \eqref{eq:u-alpha-tail}; the
factor $e^{-2|x|}$ is the reciprocal of the leading scale
$\pi e^{2|x|}$.  In particular,
\[
 \int_\R\frac{|u_\alpha(x)|^2}{W(x)}\dd x<\infty.
\]
Moreover,
\begin{equation}\label{eq:u-alpha-ode}
 Du_\alpha=-\alpha u_\alpha-iW.
\end{equation}
The same estimate shows that $u_\alpha,u_\alpha'\in L^2(\R)$ and that
$W^{-1/2}Du_\alpha\in L^2(\R)$.  Since
$\sqrt p\,g_\alpha=u_\alpha/\sqrt{c_0}$, these are exactly the two
conditions in \eqref{eq:A-domain}.
\end{proof}

Equation \eqref{eq:u-alpha-ode}, followed by integration over $\R$, gives
$\int_\R u_\alpha=-ic_0/\alpha$; here $\alpha\neq0$ because
$\Xi(0)>0$.  It follows that
\begin{equation}\label{eq:boundary-pencil}
 (A+\alpha P_0)g_\alpha=0.
\end{equation}
Applying $\BTheta$ to \eqref{eq:boundary-pencil} and using
\cref{eq:BA-I,eq:Bpsi-zero} then gives
\begin{equation}\label{eq:boundary-eigenstate}
 \BTheta g_\alpha=-\frac1\alpha g_\alpha.
\end{equation}
This is the first-order, rank-one form of the Theta boundary mechanism in
\cite{Hedenmalm2026}; the determinant argument above additionally records all
algebraic multiplicities without requiring a separate root-chain analysis.

\section{Trace identities and structural consequences}

The whole-line determinant theorem immediately extends
\cref{cor:finite-cumulant}.

\begin{theorem}[Whole-line trace--cumulant calculus]\label{thm:cumulants}
Let $p$ be Volterra-admissible, and let $\kappa_n(p)$ be its cumulants.  Then,
for every $n\geq2$,
\begin{equation}\label{eq:whole-cumulants}
 \boxed{
 \tr(\Vp^n)=
 \frac{(-1)^{n+1}i^n}{(n-1)!}\,\kappa_n(p).}
\end{equation}
\end{theorem}

For the even Riemann density, all odd cumulants vanish.  In particular,
\begin{equation}\label{eq:B2-variance}
 \tr(\BTheta^2)
 =\operatorname{Var}_p(X)
 =-\frac{\Xi''(0)}{\Xi(0)}.
\end{equation}
The same identity follows directly from the kernel.  If $y<x$, then
\[
 b_p(x,y)b_p(y,x)=F(y)(1-F(x)),
\]
and integration over the two half-triangles gives the variance.  This is an
independent normalization check: no extra factor $e^{az^2}$ can be inserted
in \eqref{eq:main-det} without changing \eqref{eq:B2-variance}.

There is also a useful conceptual separation.  The Volterra indicator in
\eqref{eq:whole-kernel} retains order; the CDF term retains the scalar endpoint
memory removed by centering; only the final regularized determinant becomes a
scalar.  Thus the construction provides a general method for realizing an
entire characteristic function as the scalar readout of an ordered compact
channel.  The Riemann case is distinguished by the arithmetic Theta density,
not by the abstract determinant mechanism alone.

The construction is faithful modulo translation in a precise sense.  For two
densities to which the finite-interval theorem (or its admissible whole-line
extension) applies, equality of their modified determinant functions is
equivalent to equality of their centered probability laws.  Indeed,
\cref{thm:finite-det,rem:nonzero-mean} identify the determinants with the
characteristic functions of $X-\mathbb E X$, and characteristic functions
determine probability measures uniquely.  Thus probability-CDF centering
defines a canonical operator encoding of a law modulo its first cumulant;
\cref{thm:cumulants} recovers every higher cumulant as a trace.

\section{Concluding remarks and the spectral-reality problem}

The results above close five operator-theoretic points unconditionally:
\begin{enumerate}[label=(\arabic*),leftmargin=2.2em]
 \item the CDF correction is uniquely forced by the local projection and
 endpoint requirements;
 \item $\BTheta$ is a well-defined Hilbert--Schmidt operator on the full real
 line, with an explicit bound;
 \item its finite centered restrictions converge in Hilbert--Schmidt norm;
 \item its regularized determinant is exactly $\Xi(z)/\Xi(0)$, including
 normalization and algebraic multiplicity;
 \item it is the Moore--Penrose inverse of a concrete closed first-order
 Theta operator, admits a Witten-type supersymmetric completion of index
 $-1$, and obeys exact PT and parity symmetries.
\end{enumerate}

None of these statements implies that the spectrum is real.  In the present
language, the remaining RH problem is exactly the following.

\begin{problem}[Theta spectral-reality problem]\label{prob:reality}
Prove, using a property specific to the arithmetic density $W$, that every
nonzero eigenvalue of $\BTheta$ is real.
\end{problem}

Self-adjointness of $\BTheta$ in the ambient $L^2$ inner product is false in
general, and PT symmetry by itself is insufficient.  One possible direction
is to seek a bounded, positive, invertible metric $G$ satisfying
$G\BTheta=\BTheta^*G$, as in pseudo-Hermitian quantum mechanics
\cite{Mostafazadeh2002,Mostafazadeh2002III}.  Such a metric would make
$\BTheta$ boundedly similar to a self-adjoint compact operator and hence force
spectral reality.  It is potentially stronger than RH alone: it would also
exclude nontrivial Jordan chains and impose Riesz-basis control on the
eigenvectors.  Its existence should therefore not be treated as a free
consequence of PT symmetry or of the rank-one construction.  Another possible
direction is to exploit additional total-positivity or variation-diminishing
structure of the specific Theta kernel.  The present paper does not assume
either property.

The value of Problem~\ref{prob:reality} is that it is now separated from all
regularization and closure issues.  Any proposed RH argument using this model
must act on the explicit kernel \eqref{eq:whole-kernel} and prove spectral
reality without defining the missing structure from the zeros themselves.

\appendix

\section{Classical derivation of the Riemann--Theta kernel}
\label{sec:appendix-kernel}

This appendix derives \eqref{eq:W-def} and \eqref{eq:fourier-xi} directly
from the classical theta transformation and Mellin formula.  Put
\begin{equation}\label{eq:appendix-theta-psi}
 \vartheta(t)=\sum_{n\in\mathbb Z}e^{-\pi n^2t},
 \qquad
 \psi(t)=\frac{\vartheta(t)-1}{2}
 =\sum_{n=1}^\infty e^{-\pi n^2t},
 \qquad t>0.
\end{equation}
Poisson summation gives Jacobi's inversion formula
\begin{equation}\label{eq:appendix-jacobi}
 \vartheta(1/t)=t^{1/2}\vartheta(t).
\end{equation}
For $\Re s>1$, termwise integration gives the classical Mellin identity
\begin{equation}\label{eq:appendix-mellin-zeta}
 \pi^{-s/2}\Gamma(s/2)\zeta(s)
 =\int_0^\infty \psi(t)t^{s/2}\frac{\dd t}{t}.
\end{equation}

Define
\begin{align}
 \omega(t)
 &:=\frac12\bigl(2t^2\vartheta''(t)+3t\vartheta'(t)\bigr)\notag\\
 &=\sum_{n=1}^\infty
 \bigl(2\pi^2n^4t^2-3\pi n^2t\bigr)e^{-\pi n^2t}.
 \label{eq:appendix-omega}
\end{align}
The differentiated theta series converges locally uniformly on $(0,\infty)$,
so $\omega$ is smooth.  Differentiating \eqref{eq:appendix-jacobi} once and
twice yields
\begin{align}
 \vartheta'(1/t)
 &=-\frac12t^{3/2}\vartheta(t)-t^{5/2}\vartheta'(t),\notag\\
 \vartheta''(1/t)
 &=\frac34t^{5/2}\vartheta(t)+3t^{7/2}\vartheta'(t)
   +t^{9/2}\vartheta''(t).
 \label{eq:appendix-theta-derivatives}
\end{align}
Substitution in \eqref{eq:appendix-omega} gives the exact transformation law
\begin{equation}\label{eq:appendix-omega-inversion}
 \omega(1/t)=t^{1/2}\omega(t).
\end{equation}

Let $a=s/2$.  Since
$\omega(t)=2t^2\psi''(t)+3t\psi'(t)$, two integrations by parts in
\eqref{eq:appendix-mellin-zeta} give, initially for $\Re s>1$,
\begin{align}
 \int_0^\infty\omega(t)t^{a-1}\dd t
 &=(1-2a)\int_0^\infty t^a\psi'(t)\dd t\notag\\
 &=a(2a-1)\int_0^\infty\psi(t)t^{a-1}\dd t\notag\\
 &=\frac12s(s-1)\pi^{-s/2}\Gamma(s/2)\zeta(s)
 =\xi(s).
 \label{eq:appendix-xi-mellin}
\end{align}
The boundary terms vanish at infinity by exponential decay.  At the origin,
\eqref{eq:appendix-jacobi} gives
\[
 \psi(t)=\frac12(t^{-1/2}-1)
 +O\!\left(t^{-1/2}e^{-\pi/t}\right),
 \qquad t\downarrow0,
\]
which also makes the boundary terms vanish when $\Re s>1$.
Equation \eqref{eq:appendix-omega-inversion} shows that the left side of
\eqref{eq:appendix-xi-mellin} converges locally uniformly for every
$s\in\C$; hence \eqref{eq:appendix-xi-mellin} holds on $\C$ by analytic
continuation.

Finally define, for every $x\in\R$,
\begin{equation}\label{eq:appendix-W-omega}
 W(x)=2e^{x/2}\omega(e^{2x}).
\end{equation}
For $x\geq0$, inserting \eqref{eq:appendix-omega} into
\eqref{eq:appendix-W-omega} is exactly the series \eqref{eq:W-def}.
Moreover, \eqref{eq:appendix-omega-inversion} gives
\[
 W(-x)=2e^{-x/2}\omega(e^{-2x})
 =2e^{x/2}\omega(e^{2x})=W(x),
\]
so the reflected definition in \eqref{eq:W-def} is the restriction of one
smooth even function.  For $t\geq1$, every summand of
\eqref{eq:appendix-omega} is positive because $2\pi n^2t-3>0$;
\eqref{eq:appendix-omega-inversion} extends positivity to $0<t<1$.
The same series gives double-exponential decay of $W(x)$ as $x\to+\infty$,
and evenness gives it as $x\to-\infty$.  Thus $W$ has exponential moments of
every order.  In \eqref{eq:appendix-xi-mellin}, set
$s=\tfrac12+iz$ and substitute $t=e^{2x}$.  Then
\begin{align*}
 \Xi(z)
 &=\int_0^\infty\omega(t)t^{(1/2+iz)/2}\frac{\dd t}{t}\\
 &=\int_{-\infty}^{\infty}
 2e^{x/2}\omega(e^{2x})e^{izx}\dd x
 =\int_{\R}W(x)e^{izx}\dd x,
\end{align*}
which proves \eqref{eq:fourier-xi} with the normalization used throughout.

\section{Details of the Theta truncation bounds}\label{sec:appendix-tail}

For completeness we record the two estimates used in
\cref{thm:theta-admissible}.  From \eqref{eq:w1-def}, the substitution
$t=\pi e^{2x}$ gives
\begin{align}
 \int_R^\infty w_1(x)\dd x
 &=\pi^{-1/4}\int_{Q_R}^\infty
 t^{1/4}(2t-3)e^{-t}\dd t\notag\\
 &\leq\frac{2\pi^{-1/4}Q_R^{5/4}e^{-Q_R}}
 {1-5/(4Q_R)}.
 \label{eq:appendix-q-tail}
\end{align}
Together with $W\leq(1+\Delta)w_1$, this proves
\eqref{eq:qR-bound}.

For the reciprocal integral, $W\geq w_1$ gives
\begin{align}
 \int_{-R}^R\frac{\dd x}{p(x)}
 &\leq2c_0\int_0^R\frac{\dd x}{w_1(x)}\notag\\
 &=\frac{c_0\pi^{1/4}}2
 \int_\pi^{Q_R}\frac{e^t}{t^{9/4}(2t-3)}\dd t\notag\\
 &\leq\frac{c_0}{2\pi^2(2\pi-3)}e^{Q_R},
 \label{eq:appendix-reciprocal}
\end{align}
which is \eqref{eq:reciprocal-bound}.  Substitution in
\eqref{eq:interior-error} gives \eqref{eq:explicit-interior-bound}.  These
fixed-constant inequalities are the analytic content behind the full-line
closure; numerical tail fitting is not used.

\section{Continuity of the regularized determinant}

We recall the standard estimate that justifies the limiting passage.  For
$S,T\in\HS$,
\begin{equation}\label{eq:det2-continuity}
 |\detTwo(I+S)-\detTwo(I+T)|
 \leq \|S-T\|_{\HS}
 \exp\!\left(C(1+\|S\|_{\HS}+\|T\|_{\HS})^2\right)
\end{equation}
for an absolute constant $C$; see \cite[Chapter 9]{Simon2005}.  On a compact
$z$-set, apply this estimate to $S=z\Vp^{(R)}$ and $T=z\Vp$.  The norms remain
bounded by \cref{prop:HS-truncation}, proving locally uniform convergence.

\section{Code and reproducibility}

The complete reproducibility bundle is supplied with this manuscript as the
ancillary archive
\texttt{btheta-volterra-fredholm-reproducibility.zip}.  It contains the
\LaTeX{} source and bibliography, the verification program
\texttt{verify\_btheta.py}, a pinned dependency file, and a README giving the
exact commands.  The program performs two independent numerical regression
tests: it checks the Fourier normalization in \cref{eq:fourier-xi} at high
precision, and it applies midpoint Nystrom discretization to the
finite-interval identity in \cref{thm:finite-det}; compare the general
numerical framework in \cite{Bornemann2010}.  No external data are required.
The source, verification program, and compiled manuscript are publicly
available at
\url{https://github.com/AlexxxxxLiu/btheta-volterra-fredholm}.
The ancillary archive may also be distributed directly as supplementary
material.  These computations are not used in any proof.  The analytic
argument depends only on the displayed inequalities and standard trace-ideal
continuity.

\enlargethispage{4\baselineskip}
\bibliographystyle{amsplain}
\bibliography{references}

\end{document}
